\documentclass[letterpaper,twocolumn,10pt]{article}
\usepackage{hyperref}
\usepackage{usenix2019_v3_xetex}
\usepackage[T1]{fontenc}
\usepackage{amsmath,amssymb,graphicx}
\hypersetup{hidelinks,pdftitle={Integrated architecture: typesetting proof},pdfauthor={}}
\makeatletter
\newcommand{\captiontypefigure}{\def\@captype{figure}}
\makeatother
\begin{document}
\raggedend
\twocolumn[{
\noindent\textbf{Integrated architecture: author-review proof}\par
\vspace{8pt}
\centering
\includegraphics[width=\textwidth]{figures/system-architecture-integrated.pdf}
\vspace{3pt}
\captiontypefigure
\caption{Shared state, packed access, and per-object storage in one hierarchy.
Each workspace resolves private current state $\Delta_i$ over the same pinned
snapshot $S$. The aligned range maps retain explicit lower references and
point private changes to separate payload objects $U_i$ outside $S$.
Within $S$, a representative packed record combines a name, attributes and
a locator to a frame in $O_1$; metadata pages and payload objects need not
be the same object. Grey projections connect $O_1$ and $O_2$ to their own
stored representations, not to additional read-time processing stages.
Their equal-width strips illustrate different supported fixed encoding
blocks $B_1$ and $B_2$. The enlarged full-stripe shard inset has
$s=B_1/d=2q_1$, where $q$ is the verification quantum; it is distinct from
a packed frame. At upload, a configured client policy may request $B$ using the
finalized object boundary. RustFS admits and persists the effective $(B,q)$;
reads and repair use the stored layout. This is not an automatic-policy
outcome. Dashed hint wiring still needs Entity/producer integration, and
smaller-$q$ writes require eligible new objects and both default-off gates
enabled. Upper reads verify whole objects before slicing; $\Delta_i$ remains
in the mutable metadata service. Other access and lifecycle paths are omitted.}

\vspace{14pt}
}]
\section{Shared representation}
Workspace $i$ resolves $S$ and its private current state $\Delta_i$.
The figure retains the explicit reference structure: unchanged intervals
refer to $S$, and private payload objects remain separate. Within the shared
snapshot, packed records organize names, common attributes and locations
for subsequent access. The locator shown is one representative access path.

\newpage
\section{Persistent physical interpretation}
The same object identities continue into the storage layer. Their vertical
alignment expresses logical and physical representations, rather than a
mandatory sequence of read operations. The upload branch is separate:
RustFS admits the requested coding block and persists compatible geometry.
Reads and repair use the saved layout, including the verification quantum.
The differing block patterns demonstrate per-object configurability.
\end{document}
